\subsection*{1.2. The geometry of Hilbert spaces}
\hfill 9

Note that $f_n$ will be as smooth as $u_n$, hence the title smoothing lemma.
The same idea is used to approximate noncontinuous functions by smooth
ones (of course the convergence will no longer be uniform in this case).

\noindent\textbf{Corollary 1.4.} $\mathcal{C}(I)$ is separable.

The same is true for $l^1(\mathbb{N})$, but not for $l^\infty(\mathbb{N})$ (Problem 1.4)!

\noindent\textbf{Problem 1.1.} Show that $|||f|| - ||g||| \leq ||f - g||$.

\noindent\textbf{Problem 1.2.} Show that the norm, vector addition, and multiplication by
scalars are continuous. That is, if $f_n \to f$, $g_n \to g$, and $\lambda_n \to \lambda$ then
$\|f_n\| \to \|f\|$, $f_n + g_n \to f + g$, and $\lambda_n g_n \to \lambda g$.

\noindent\textbf{Problem 1.3.} Show that $l^\infty(\mathbb{N})$ is a Banach space.

\noindent\textbf{Problem 1.4.} Show that $l^1(\mathbb{N})$ is separable. Show that $l^\infty(\mathbb{N})$ is not sep-
arable (Hint: Consider sequences which take only the value one and zero.
How many are there? What is the distance between two such sequences?).

\subsection*{1.2. The geometry of Hilbert spaces}

So it looks like $\mathcal{C}(I)$ has all the properties we want. However, there is
still one thing missing: How should we define orthogonality in $\mathcal{C}(I)$? In
Euclidean space, two vectors are called orthogonal if their scalar product
vanishes, so we would need a scalar product:

Suppose $\mathcal{H}$ is a vector space. A map $\langle \cdot, \cdot \rangle: \mathcal{H} \times \mathcal{H} \to \mathbb{C}$ is called skew
linear form if it is conjugate linear in the first and linear in the second
argument, that is,
\begin{align*}
\langle \lambda_1 f_1 + \lambda_2 f_2, g \rangle &= \lambda_1^* \langle f_1, g \rangle + \lambda_2^* \langle f_2, g \rangle \\
\langle f, \lambda_1 g_1 + \lambda_2 g_2 \rangle &= \lambda_1 \langle f, g_1 \rangle + \lambda_2 \langle f, g_2 \rangle, \quad \lambda_1, \lambda_2 \in \mathbb{C},
\end{align*}
(1.17)
where $*$ denotes complex conjugation. A skew linear form satisfying the
requirements
\begin{enumerate}
    \item $\langle f, f \rangle > 0$ for $f \neq 0$.
    \item $\langle f, g \rangle = \langle g, f \rangle^*$
\end{enumerate}
is called inner product or scalar product. Associated with every scalar
product is a norm
\[
\|f\| = \sqrt{\langle f, f \rangle}.
\tag{1.18}
\]
The pair $(\mathcal{H}, \langle \cdot, \cdot \rangle)$ is called inner product space. If $\mathcal{H}$ is complete it is
called a Hilbert space.

\noindent\textbf{Example.} Clearly $\mathbb{C}^n$ with the usual scalar product
\[
(a,b) = \sum_{j=1}^n a_j \bar{b}_j
\tag{1.19}
\]